% Source: user-provided adjacent-LCA virtual-tree construction.
% Supplement: https://oi-wiki.org/graph/virtual-tree/
\subsection{虚树：保留关键点与分叉关系}
\label{tree:virtual-tree}
\AlgorithmInfo
    {在固定大树上，把一组关键点相关的祖先关系压缩成小树，供多次稀疏查询使用。}
    {预处理 $O(n\log n)$；每轮输入 $k$ 个点，构造 $O(k\log(k+1)+k\log n+k_{prev})$，总空间 $O(n\log n)$。}
    {先复制 TreeLCA；原图为连通树，固定根；允许重复、空关键点集，构造时总是加入根。}
\subsubsection*{原理与构造}
只保留关键点时会丢失分叉，因此还要保留它们的 LCA。将关键点与根按 DFS 序排序，补入每对相邻点的 LCA，再排序去重，就包含了需要的全部分叉点：若关键点分布在某个 LCA 的不同儿子子树内，DFS 连续性保证至少有一对相邻关键点跨过这个分界。

设最终序列为 $a_1,\ldots,a_m$，对每对相邻点连接
\[
    LCA(a_i,a_{i+1})\longrightarrow a_{i+1}.
\]
补齐 LCA 后，该祖先就是 $a_{i+1}$ 在保留节点中的最近祖先。每个非根点恰有一条入边，保持原树祖先关系，而把中间没有保留节点的长链压成一条边。非空关键点集去重后有 $q$ 个点时，加入根后的虚树至多 $2q$ 个点；空集合只保留根。

\subsubsection*{压缩边上保留什么}
虚树边 $p\to v$ 代表原树一段路径，权值按问题定义：距离和为 $dis_v-dis_p$，边数为 $dep_v-dep_p$，路径最小边权则需要相应的路径聚合，不能拿两个前缀最小值相减。下一节的断边示例直接使用原树根路径最小值 \InlineCode{minn[v]}，并非把它当作通用压缩边权。

\subsubsection*{使用教程}
\InlineCode{init(n)} 后用 \InlineCode{add(u,v,w)} 加无向边，每边一次、默认权值 1；\InlineCode{prepare(root)} 预处理原树，省略根取 1。之后 \InlineCode{build(keys)} 接收无占位的关键点列表，允许重复与空集，构建本轮虚树；\InlineCode{nodes} 是本轮保留点的 DFS 序列表，\InlineCode{vt[u]} 是原编号下的虚树儿子，\InlineCode{key[u]} 标记实际输入的关键点。根被自动加入不代表它自动成为关键点。
\begin{lstlisting}[style=mtycpp,numbers=none]
VirtualTree g;
g.init(5);
g.add(1,2,3),g.add(1,3,8),g.add(2,4,5),g.add(2,5,2);
g.prepare(1);
g.build(vector<int>{3,4,4});
auto nodes=g.nodes;
for(auto u:g.nodes) for(auto v:g.vt[u]){
	long long length=g.dis[v]-g.dis[u];
}
g.build(vector<int>{});
\end{lstlisting}
每轮只清理上次 \InlineCode{nodes} 中的邻接表与标记，避免按原树 $n$ 清空；复杂度中的 $k_{prev}$ 就是该清理量。\InlineCode{clear()} 可手动丢弃本轮虚树，原树和倍增保留。\InlineCode{prepare} 会重新预处理原树，\InlineCode{build} 专门用于关键点查询，两者不要混淆。
\needspace{7\baselineskip}
\subsubsection*{模板代码}
\CodeFile{code/trees/virtual_tree.cpp}
